\chapter{Mechanical response and the resilient modulus of biochar-treated soil under cyclic loading}
\label{chap:four}



\section{Test apparatus}
\subsection{Unsaturated soil triaxial testing apparatus}

For the triaxial tests, specimens were prepared in a three-part split mould, with a inner diameter of 50 mm and a height of 100 mm. \par


In Section \ref{sec:ch3oedometerapparatus}, the specimens were tested under a $K_{0}$ condition, while the actual stress state is unknown. Therefore, the effect of stress state on the \ac{SWRC} was further investigated using a suction-controlled triaxial apparatus (GDS Instruments, UK), as shown in Figure \ref{fig:ch4triaxial}. Matric suction was controlled by axis translation through a \ac{HAEV} ceramic disc with an air-entry value of 1,000 kPa. The apparatus was equipped with an Advanced Air Pressure/Volume Controller, which can control the pore-air pressure while simultaneously measuring changes in air volume. Therefore, the total volume change ($\Delta V_{\mathrm{t}}$) of the specimen could be determined as the sum of the changes in pore-air volume ($\Delta V_{\mathrm{a}}$) and pore-water volume ($\Delta V_{\mathrm{w}}$):
\begin{equation}
    \Delta V_{\mathrm{t}} = \Delta V_{\mathrm{a}} + \Delta V_{\mathrm{w}}
\end{equation}
According to the ideal gas law, the pore-air pressure was maintained constant to minimise errors in measuring $\Delta V_{\mathrm{a}}$. Therefore, the net stress state ($\bm{\sigma}_{\mathrm{net}}$) and matric suction ($s$) were adjusted by varying the confining pressure ($p_{\mathrm{c}}$) and the pore-water pressure, respectively. The net stress state of the specimen is calculated via:
\begin{equation}
    \bm{\sigma}_{\mathrm{net}} = -(p_{\mathrm{a}} - p_{\mathrm{c}})\cdot \mathbf{I}
\end{equation}
where $\mathbf{I}$ is an identity matrix.\par
\begin{figure}[htbp]
    \centering
    \includegraphics[width=0.76\linewidth]{../figures/Ch4/ch4triaxial.png}
    \caption{Unsaturated soil triaxial testing apparatus}
    \label{fig:ch4triaxial}
\end{figure}
